% Supplementary section: real-video observational audit on PA-HMDB51.
% Drop into the supplementary after the zero-shot section. Numbers from cluster
% jobs 271567 (25-class screen) and 271696 (dribble study); analysis scripts:
% ActionBiasBench/scripts/eval_pahmdb51_zero_shot.py and eval_dribble_scaffold.py.
% Add to the bibliography: Wu et al., "Privacy-Preserving Deep Action Recognition:
% An Adversarial Learning Framework and A New Dataset", IEEE TPAMI 2022 -> \cite{pahmdb51}.

\section{Real-Video Observational Audit on PA-HMDB51}
\label{sec:supp_pahmdb51}

The main paper argues that observational audits of recorded video generally
cannot determine which visual factor caused a performance difference. Here we
run exactly such an audit and show this failure mode concretely on real data.
We evaluate the same six Kinetics-400-pretrained backbones from
\cref{sec:results} zero-shot, with their original 400-way classification
heads, on real HMDB51 clips carrying per-video skin-tone annotations from
PA-HMDB51~\cite{pahmdb51}. We keep the 25 HMDB51 classes with an unambiguous
Kinetics-400 counterpart (e.g., \textit{dribble}~$\rightarrow$~\textit{dribbling
basketball}; the four sword-related classes are pooled onto \textit{sword
fighting}), yielding 242 annotated clips. Pooling all classes and models, the
accuracy gap between clips labeled \textit{white} and all other skin-tone
labels is small and not significant ($0.79$ vs.\ $0.74$ restricted to the 25
mapped labels; Fisher exact $p=0.058$, and this value is optimistic since it
treats the six models' predictions on the same clip as independent). The
skin-tone label is also strongly confounded with the action class itself:
several classes are annotated almost exclusively white, whereas
\textit{dribble} contains no white-labeled clips at all, so pooled gaps
conflate class difficulty with skin tone.

To increase power where the demographic contrast is largest, we annotated the
skin tone of all 145 \textit{dribble} clips in HMDB51 (the original PA-HMDB51
covers six): 75 white, 39 black, 30 mixed, 1 unidentifiable. A naive
clip-level comparison produces an apparently decisive result: white-labeled
clips are recognized far more accurately than the rest, with
$p\approx8\times10^{-8}$ (\cref{tab:pahmdb51_dribble}). However, HMDB51 cuts
multiple clips from the same source video, and the 145 dribble clips originate
from only 29 source videos, each with a single performer, court, and camera
setup; skin tone is constant within every source. Treating the source video as
the unit of analysis --- the actually independent sample --- the same
comparison collapses to chance ($p=0.49$; 16 white vs.\ 7 black sources), and
the residual descriptive gap is partly driven by two atypical movie-sourced
clips (one in each group), not by tone.

\begin{table}[t]
\centering
\caption{Zero-shot accuracy on the 145 real HMDB51 \textit{dribble} clips by
annotated skin tone, pooled over the six Kinetics-400 backbones. \emph{Exact}
scores a prediction correct only for \textit{dribbling basketball};
\emph{family} accepts any of the four K400 basketball classes. The clip-level
test treats each (clip, model) pair as independent although the clips stem
from only 29 source videos; the source-level test uses one accuracy value per
source video (Mann--Whitney $U$, white vs.\ black sources). The
``significant'' clip-level disparity disappears once performer/background
pseudo-replication is removed.}
\label{tab:pahmdb51_dribble}
\begin{tabular}{llccc}
\toprule
Unit of analysis & Metric & White & Non-white & $p$ \\
\midrule
Clip $\times$ model ($n{=}450$ vs.\ $414$) & exact  & $0.82$ & $0.66$ & $8.1\times10^{-8}$ \\
Clip $\times$ model ($n{=}450$ vs.\ $414$) & family & $0.96$ & $0.88$ & $1.9\times10^{-5}$ \\
Source video ($n{=}16$ vs.\ $7$)           & exact  & $0.67$ & $0.61$ & $0.49$ \\
Source video ($n{=}16$ vs.\ $7$)           & family & $0.91$ & $0.77$ & $0.23$ \\
\bottomrule
\end{tabular}
\end{table}

We emphasize what this does and does not show. It does not demonstrate the
absence of skin-tone bias on real video: with only seven independent
black-performer sources, even a double-digit accuracy gap would remain
undetectable. What it does demonstrate is that a real-video audit of typical
benchmark size can manufacture a highly ``significant'' skin-tone disparity
that is, in fact, performer- and background-driven pseudo-replication. This is
precisely the confounding that recorded data cannot resolve and that motivates
the controlled counterfactual design of the main paper, where the same motion,
scene, and camera are held fixed while only skin tone changes.
